% Options for packages loaded elsewhere
\PassOptionsToPackage{unicode}{hyperref}
\PassOptionsToPackage{hyphens}{url}
\documentclass[
]{article}
\usepackage{xcolor}
\usepackage[margin=1in]{geometry}
\usepackage{amsmath,amssymb}
\setcounter{secnumdepth}{-\maxdimen} % remove section numbering
\usepackage{iftex}
\ifPDFTeX
  \usepackage[T1]{fontenc}
  \usepackage[utf8]{inputenc}
  \usepackage{textcomp} % provide euro and other symbols
\else % if luatex or xetex
  \usepackage{unicode-math} % this also loads fontspec
  \defaultfontfeatures{Scale=MatchLowercase}
  \defaultfontfeatures[\rmfamily]{Ligatures=TeX,Scale=1}
\fi
\usepackage{lmodern}
\ifPDFTeX\else
  % xetex/luatex font selection
\fi
% Use upquote if available, for straight quotes in verbatim environments
\IfFileExists{upquote.sty}{\usepackage{upquote}}{}
\IfFileExists{microtype.sty}{% use microtype if available
  \usepackage[]{microtype}
  \UseMicrotypeSet[protrusion]{basicmath} % disable protrusion for tt fonts
}{}
\makeatletter
\@ifundefined{KOMAClassName}{% if non-KOMA class
  \IfFileExists{parskip.sty}{%
    \usepackage{parskip}
  }{% else
    \setlength{\parindent}{0pt}
    \setlength{\parskip}{6pt plus 2pt minus 1pt}}
}{% if KOMA class
  \KOMAoptions{parskip=half}}
\makeatother
\usepackage{color}
\usepackage{fancyvrb}
\newcommand{\VerbBar}{|}
\newcommand{\VERB}{\Verb[commandchars=\\\{\}]}
\DefineVerbatimEnvironment{Highlighting}{Verbatim}{commandchars=\\\{\}}
% Add ',fontsize=\small' for more characters per line
\usepackage{framed}
\definecolor{shadecolor}{RGB}{248,248,248}
\newenvironment{Shaded}{\begin{snugshade}}{\end{snugshade}}
\newcommand{\AlertTok}[1]{\textcolor[rgb]{0.94,0.16,0.16}{#1}}
\newcommand{\AnnotationTok}[1]{\textcolor[rgb]{0.56,0.35,0.01}{\textbf{\textit{#1}}}}
\newcommand{\AttributeTok}[1]{\textcolor[rgb]{0.13,0.29,0.53}{#1}}
\newcommand{\BaseNTok}[1]{\textcolor[rgb]{0.00,0.00,0.81}{#1}}
\newcommand{\BuiltInTok}[1]{#1}
\newcommand{\CharTok}[1]{\textcolor[rgb]{0.31,0.60,0.02}{#1}}
\newcommand{\CommentTok}[1]{\textcolor[rgb]{0.56,0.35,0.01}{\textit{#1}}}
\newcommand{\CommentVarTok}[1]{\textcolor[rgb]{0.56,0.35,0.01}{\textbf{\textit{#1}}}}
\newcommand{\ConstantTok}[1]{\textcolor[rgb]{0.56,0.35,0.01}{#1}}
\newcommand{\ControlFlowTok}[1]{\textcolor[rgb]{0.13,0.29,0.53}{\textbf{#1}}}
\newcommand{\DataTypeTok}[1]{\textcolor[rgb]{0.13,0.29,0.53}{#1}}
\newcommand{\DecValTok}[1]{\textcolor[rgb]{0.00,0.00,0.81}{#1}}
\newcommand{\DocumentationTok}[1]{\textcolor[rgb]{0.56,0.35,0.01}{\textbf{\textit{#1}}}}
\newcommand{\ErrorTok}[1]{\textcolor[rgb]{0.64,0.00,0.00}{\textbf{#1}}}
\newcommand{\ExtensionTok}[1]{#1}
\newcommand{\FloatTok}[1]{\textcolor[rgb]{0.00,0.00,0.81}{#1}}
\newcommand{\FunctionTok}[1]{\textcolor[rgb]{0.13,0.29,0.53}{\textbf{#1}}}
\newcommand{\ImportTok}[1]{#1}
\newcommand{\InformationTok}[1]{\textcolor[rgb]{0.56,0.35,0.01}{\textbf{\textit{#1}}}}
\newcommand{\KeywordTok}[1]{\textcolor[rgb]{0.13,0.29,0.53}{\textbf{#1}}}
\newcommand{\NormalTok}[1]{#1}
\newcommand{\OperatorTok}[1]{\textcolor[rgb]{0.81,0.36,0.00}{\textbf{#1}}}
\newcommand{\OtherTok}[1]{\textcolor[rgb]{0.56,0.35,0.01}{#1}}
\newcommand{\PreprocessorTok}[1]{\textcolor[rgb]{0.56,0.35,0.01}{\textit{#1}}}
\newcommand{\RegionMarkerTok}[1]{#1}
\newcommand{\SpecialCharTok}[1]{\textcolor[rgb]{0.81,0.36,0.00}{\textbf{#1}}}
\newcommand{\SpecialStringTok}[1]{\textcolor[rgb]{0.31,0.60,0.02}{#1}}
\newcommand{\StringTok}[1]{\textcolor[rgb]{0.31,0.60,0.02}{#1}}
\newcommand{\VariableTok}[1]{\textcolor[rgb]{0.00,0.00,0.00}{#1}}
\newcommand{\VerbatimStringTok}[1]{\textcolor[rgb]{0.31,0.60,0.02}{#1}}
\newcommand{\WarningTok}[1]{\textcolor[rgb]{0.56,0.35,0.01}{\textbf{\textit{#1}}}}
\usepackage{graphicx}
\makeatletter
\newsavebox\pandoc@box
\newcommand*\pandocbounded[1]{% scales image to fit in text height/width
  \sbox\pandoc@box{#1}%
  \Gscale@div\@tempa{\textheight}{\dimexpr\ht\pandoc@box+\dp\pandoc@box\relax}%
  \Gscale@div\@tempb{\linewidth}{\wd\pandoc@box}%
  \ifdim\@tempb\p@<\@tempa\p@\let\@tempa\@tempb\fi% select the smaller of both
  \ifdim\@tempa\p@<\p@\scalebox{\@tempa}{\usebox\pandoc@box}%
  \else\usebox{\pandoc@box}%
  \fi%
}
% Set default figure placement to htbp
\def\fps@figure{htbp}
\makeatother
\setlength{\emergencystretch}{3em} % prevent overfull lines
\providecommand{\tightlist}{%
  \setlength{\itemsep}{0pt}\setlength{\parskip}{0pt}}
\usepackage{bookmark}
\IfFileExists{xurl.sty}{\usepackage{xurl}}{} % add URL line breaks if available
\urlstyle{same}
\hypersetup{
  hidelinks,
  pdfcreator={LaTeX via pandoc}}

\author{}
\date{\vspace{-2.5em}}

\begin{document}

\section{TaxaExpect}\label{taxaexpect}

Estimate spatially-explicit Bayesian priors for expected species
composition. Part of the
\href{https://github.com/DOI-USGS/TaxaID}{TaxaID} ecosystem.

A taxonomic assignment based on a match to a reference can propose an
implausible species. These are sometimes accepted as fact, but often are
screened afterwards, or prevented from occuring by reducing the
reference list to known local species.

Bayes' Theorem improves taxonomic assignment by considering, before any
match score is seen, how large a share of the detections at the sampling
location a hypothesized taxon is expected to account for. TaxaExpect
estimates these priors from occurrence records. This is typically GBIF
(Global Biodiversity Information Facility; GBIF Secretariat, Copenhagen,
Denmark) records fetched via TaxaFetch, or user-supplied data. Although
occurrence data are often sparse and biased, they are usually sufficient
to distinguish among taxa with similar match scores but very different
geographic ranges. TaxaExpect's priors are compositional shares -- the
expected relative share of a species at a site, not an occurrence
probability -- estimated directly from occurrence data rather than from
an occurrence/occupancy model.

TaxaExpect uses a \textbf{site-centered distance kernel}: each species'
prior is its kernel-weighted share of related, nearby occurrence
records, computed directly at the query site. Weight decays smoothly
with distance (and, optionally, with a covariate such as depth), so
spatial borrowing degrades continuously rather than switching on or off
at an arbitrary cell boundary, and even lightly-sampled sites get
informative priors by shrinking toward the surrounding region's
composition. The kernel bandwidth is chosen from data via
leave-one-block-out composition prediction
(\texttt{calibrate\_kernel\_bandwidth()}), not guessed, and the
resulting prior field can be mapped continuously over the study area
with \texttt{plot\_theta\_surface()}. For taxa never reported in the
study area, TaxaExpect prices a ``dark diversity'' prior from a
Good-Turing/Chao- anchored presence-distance curve.

The resulting spatial priors are helpful in reducing false positives
from ecologically implausible assignments and can break ties between
similar-scoring species, rescuing species-level resolution that would
otherwise be lost to defensive upranking.

\subsection{Overview}\label{overview}

TaxaExpect generates theta priors for taxonomic assignment from
occurrence data. Theta is compositional: the expected relative share of
a species at a site, P(a random legitimate detection = species X). It is
not an occurrence probability, and not occupancy -- see \emph{Shared
detection effort} below for the formal definition. The current (kernel)
pathway estimates expected species composition directly at a site via
distance-weighted occurrence sharing, incorporating habitat
stratification and, optionally, a covariate such as depth.

Priors are organized into three branches (\texttt{prior\_branch}):

\begin{itemize}
\tightlist
\item
  \textbf{\texttt{resident\_observed}} -- species with kernel-weighted
  local occurrence evidence (direct estimate)
\item
  \textbf{\texttt{resident\_undetected}} -- species plausibly present
  but not locally recorded (singleton mirrors, a Good-Turing floor, and
  named claimants priced on a shared presence-distance curve)
\item
  \textbf{\texttt{transport}} -- domestic/food/cultivar species and
  other non-resident presence hypotheses
\end{itemize}

\begin{figure}
\centering
\pandocbounded{\includegraphics[keepaspectratio,alt={Interactive prior-field map produced by plot\_theta\_surface(). The surface is the estimator evaluated continuously across a lattice, not one value per grid cell, so it shows how a species' expected share changes across the region rather than only at the sampling site. Here theta for kelp surfperch (Brachyistius frenatus) rises along the Santa Barbara Channel and the Los Angeles coastline, and the hollow circle marks the focal site. Passing several taxa to taxon adds the layer toggle at top right, so priors can be compared species by species. The surface is clipped to a supplied search polygon via mask, which is why it stops at the coast instead of painting inland.}]{man/figures/theta_surface.png}}
\caption{Interactive prior-field map produced by
\protect\texttt{plot\_theta\_surface()}. The surface is the estimator
evaluated continuously across a lattice, not one value per grid cell, so
it shows how a species' expected share changes across the region rather
than only at the sampling site. Here theta for kelp surfperch
(\emph{Brachyistius frenatus}) rises along the Santa Barbara Channel and
the Los Angeles coastline, and the hollow circle marks the focal site.
Passing several taxa to \protect\texttt{taxon} adds the layer toggle at
top right, so priors can be compared species by species. The surface is
clipped to a supplied search polygon via \protect\texttt{mask}, which is
why it stops at the coast instead of painting inland.}
\end{figure}

\subsection{Assumptions}\label{assumptions}

\subsubsection{Shared detection effort}\label{shared-detection-effort}

TaxaExpect models \emph{relative} abundance: the probability that a
randomly selected observation at a site belongs to a given taxon. The
denominator \texttt{n\_total\_at\_site} is the total count of all
observations at a site and serves as the shared effort measure for every
taxon in the model.

\textbf{All taxa in a particular model should be detected through a
similar sampling process.} Combining taxa collected by incommensurable
methods corrupts the shared denominator: a phytoplankton cell count and
a bird point-count sighting are not equivalent detection events. Mixing
them implies that plankton-sampling effort informs expected bird
relative abundance, which is not true. The two surveys represent
independent detection processes.

\emph{Problematic mixing (do not combine in a single model):}

\begin{itemize}
\tightlist
\item
  Phytoplankton cell counts + bird point counts
\item
  eDNA reads from different gene markers (e.g.~12S fish reads combined
  with COI invertebrate reads; amplification efficiency differs between
  markers, so read counts are not on a shared effort scale)
\item
  Arthropods from pitfall traps + arthropods from Malaise traps
\end{itemize}

\emph{Valid pooling:}

\begin{itemize}
\tightlist
\item
  All fish species from the same eDNA marker on the same filter
\item
  All bird species detected during the same standardized point count
\item
  All macroinvertebrate taxa from the same kick-net sample
\end{itemize}

\textbf{This grouping is not automatic, and there is no automated way to
make it automatic.} \texttt{estimate\_kernel\_priors()} has no way to
tell, on its own, which taxa share a detection process -- if you leave
\texttt{sampling\_group\_col} unset, it silently pools every taxon into
one shared denominator regardless of detection process, with no warning
and no error. Supplying a correct grouping is your responsibility, and
it must be done by hand, from real knowledge of how each taxon was
detected -- there is no way to infer ``comparable detection method''
from taxonomy or occurrence data alone (two co-occurring taxa sampled by
different gear look identical in an occurrence table). Automatic,
data-driven classification of detection process does not work reliably
for this: tested against a real, full-scale expert- classified dataset,
it answered a different question than \texttt{sampling\_group} is meant
to answer (whether a group clears a per-site record-count floor, not
whether it shares a detection process) -- it fragmented a single real
detection process into dozens of spurious groups, and separately
conflated two taxa an expert had deliberately kept apart (a genuine
marine survey target and likely airborne contamination) purely because
neither cleared the floor on its own. An LLM guess at the grouping is
not a substitute for real methodological knowledge either, for the same
reason: the distinction lives in how the data were collected, not in
anything visible in the records themselves. Two options remain: 1)
manually choose similar taxa that are sampled in similar ways, the
approach this project's own analysts have used successfully for every
real dataset handled so far; or 2) build separate models for each survey
type and pass the appropriate priors to TaxaAssign for the relevant
taxonomic hypotheses.

Choosing a grouping by hand does not require worrying about sample size.
\texttt{estimate\_kernel\_priors()} does \textbf{not} need pre-merged,
sample-size-adequate groups for statistical adequacy -- it produces
honestly wide uncertainty for a small group on its own (confirmed on a
real 9-group expert classification, including a single-taxon group and a
3-taxon group, both of which fit cleanly with appropriately wide
\texttt{theta\_sd}). Classify at whatever granularity genuinely reflects
distinct detection methods, and let the estimator's own uncertainty
reflect how much data backs each group -- there is no separate merging
step to get right.

\subsection{Installation}\label{installation}

\begin{Shaded}
\begin{Highlighting}[]
\CommentTok{\# Requires TaxaTools (foundation package)}
\NormalTok{devtools}\SpecialCharTok{::}\FunctionTok{install}\NormalTok{(}\StringTok{"path/to/TaxaTools"}\NormalTok{)}
\NormalTok{devtools}\SpecialCharTok{::}\FunctionTok{install}\NormalTok{(}\StringTok{"path/to/TaxaExpect"}\NormalTok{)}
\end{Highlighting}
\end{Shaded}

\subsection{Quick Start}\label{quick-start}

\begin{Shaded}
\begin{Highlighting}[]
\FunctionTok{library}\NormalTok{(TaxaExpect)}

\CommentTok{\# 1. Calibrate the kernel bandwidth (and regional back{-}off m) by}
\CommentTok{\#    leave{-}one{-}block{-}out composition prediction {-}{-} never hand{-}set.}
\NormalTok{calib }\OtherTok{\textless{}{-}} \FunctionTok{calibrate\_kernel\_bandwidth}\NormalTok{(}
  \AttributeTok{occurrence\_data =}\NormalTok{ occurrences,   }\CommentTok{\# habitat{-}labelled occurrence records}
  \AttributeTok{site\_habitat    =} \StringTok{"Marine"}\NormalTok{,}
  \AttributeTok{lambda\_grid     =} \FunctionTok{c}\NormalTok{(}\DecValTok{10}\NormalTok{, }\DecValTok{25}\NormalTok{, }\DecValTok{50}\NormalTok{, }\DecValTok{100}\NormalTok{, }\DecValTok{200}\NormalTok{)   }\CommentTok{\# km}
\NormalTok{)}
\NormalTok{lambda\_km }\OtherTok{\textless{}{-}}\NormalTok{ calib}\SpecialCharTok{$}\NormalTok{best}\SpecialCharTok{$}\NormalTok{lambda\_km}

\CommentTok{\# 2. Estimate kernel priors at your site (resident\_observed rows)}
\NormalTok{kernel\_fit }\OtherTok{\textless{}{-}} \FunctionTok{estimate\_kernel\_priors}\NormalTok{(}
  \AttributeTok{occurrence\_data =}\NormalTok{ occurrences,}
  \AttributeTok{site\_lat        =} \FloatTok{34.45}\NormalTok{,}
  \AttributeTok{site\_lon        =} \SpecialCharTok{{-}}\FloatTok{120.47}\NormalTok{,}
  \AttributeTok{site\_habitat    =} \StringTok{"Marine"}\NormalTok{,}
  \AttributeTok{lambda\_km       =}\NormalTok{ lambda\_km}
\NormalTok{)}

\CommentTok{\# 3. Add resident\_undetected rows (singleton mirrors + Good{-}Turing floor)}
\NormalTok{undetected }\OtherTok{\textless{}{-}} \FunctionTok{generate\_undetected\_diversity}\NormalTok{(kernel\_fit)}

\CommentTok{\# 4. Assemble the prior table for TaxaAssign}
\NormalTok{priors }\OtherTok{\textless{}{-}}\NormalTok{ dplyr}\SpecialCharTok{::}\FunctionTok{bind\_rows}\NormalTok{(kernel\_fit}\SpecialCharTok{$}\NormalTok{priors, undetected)}

\CommentTok{\# 5. Explore the prior field}
\FunctionTok{plot\_theta\_surface}\NormalTok{(kernel\_fit, }\AttributeTok{occurrence\_data =}\NormalTok{ occurrences,}
                    \AttributeTok{taxon =} \StringTok{"Girella nigricans"}\NormalTok{)}
\end{Highlighting}
\end{Shaded}

\subsection{Key Functions}\label{key-functions}

\subsubsection{Kernel pathway (current,
recommended)}\label{kernel-pathway-current-recommended}

\textbf{Calibration:} - \texttt{calibrate\_kernel\_bandwidth()} --
choose the geographic bandwidth (and optional covariate bandwidth, and
the regional back-off mass \texttt{m}) by leave-one-block-out
composition prediction

\textbf{Prior estimation:} - \texttt{estimate\_kernel\_priors()} --
site-centered kernel estimation of \texttt{resident\_observed} priors
(no grid, no model fit) - \texttt{generate\_undetected\_diversity()} --
singleton-mirror and global-floor \texttt{resident\_undetected} priors
(accepts kernel or GLMM input) -
\texttt{generate\_presence\_curve\_evidence()} /
\texttt{generate\_user\_specified\_evidence()} +
\texttt{apply\_undetected\_evidence(pricing\ =\ "curve")} -- price named
unobserved claimants (regional, watch-listed, or distance-clamped) on a
shared presence-distance curve -
\texttt{condition\_evidence\_on\_habitat()} -- multiply any evidence
table's presence weights by each taxon's weight for the site habitat
(from the cached LLM habitat lookup), floored at the zero-evidence
clamp, so evidence obeys the same habitat stratification as the resident
priors while habitat bleed survives in proportion -
\texttt{generate\_domestic\_food\_priors()} -- \texttt{transport}-branch
priors for domestic, food, and cultivar species

\textbf{Diagnostics and reporting:} - \texttt{plot\_theta\_surface()} --
continuous prior-field map, evaluating the estimator on a lattice via
FFT - \texttt{kernel\_budget\_sensitivity()} -- reports how the
Good-Turing budget behind \texttt{theta\_present} moves across counting
radius/bandwidth choices - \texttt{report\_priors()} -- generate report
section for \texttt{assemble\_report()}

\subsection{Statistical Methods}\label{statistical-methods}

TaxaExpect's current pathway prices each species' prior as its
kernel-weighted share of occurrence records in the focal-habitat
stratum, shrunk toward the regional composition by \texttt{m}
pseudo-records (a Dirichlet back-off):

\begin{verbatim}
theta_i = (c_i * s + m * p_i) / (n_eff + m)
\end{verbatim}

where \texttt{c\_i} is species \emph{i}'s summed distance-weighted
record count (\texttt{w\_r\ =\ exp(-d\_r/lambda\_km)}, optionally
multiplied by a covariate factor such as depth),
\texttt{s\ =\ n\_eff\ /\ W} is an effective-scale factor, \texttt{p\_i}
is the unweighted regional (habitat-stratified) record share, and
\texttt{n\_eff\ =\ W\^{}2\ /\ sum(w\_r\^{}2)} is the Kish (1965)
effective sample size of the weighted neighborhood.
\texttt{alpha\ =\ c\_i*s\ +\ m*p\_i} and
\texttt{beta\ =\ (n\_eff\ +\ m)\ -\ alpha} are the row's Beta parameters
directly -- no delta-method back-transformation, no phi cap/floor, no
Jeffreys fallback are needed, since \texttt{alpha}/\texttt{beta} are
built from non-negative counts and pseudo-counts by construction and
cannot produce the boundary pathologies a link-scale model can.
\texttt{lambda\_km} (and, if used, a covariate bandwidth and \texttt{m})
is chosen from data by \texttt{calibrate\_kernel\_bandwidth()}'s
leave-one-block-out composition prediction -- never hand-set.

For undetected species, a Good-Turing/Chao-anchored presence-distance
curve (\texttt{apply\_undetected\_evidence(pricing\ =\ "curve")}) prices
each named claimant as a two-point presence mixture,
\texttt{theta\ =\ w\ *\ theta\_present}, where \texttt{theta\_present}
is the neighborhood's own observed singleton mean and \texttt{w} comes
from the claimant's distance to its nearest occurrence record (or, for
iNaturalist range evidence, from independent verification).

\textbf{Real validation.} Leave-one-block-out testing -- holding out
blocks of occurrence records and scoring composition predictions against
them by multinomial log-loss -- found that single-cell GLMM prediction
scored \emph{worse} than ignoring space entirely, while the kernel
estimator beat both regional pooling and the single-cell predictor at
every bandwidth tested. On real Great Lakes data, switching from the
GLMM baseline to the kernel estimator raised species co-detections from
237 to 564 and precision from 0.748 to 0.868 (independently validated
against a held-out checklist).

For the full statistical derivation, assumptions, and references, see
\href{inst/TaxaExpect_supplemental_methods.md}{\texttt{inst/TaxaExpect\_supplemental\_methods.md}}.

\subsection{Part of TaxaID}\label{part-of-taxaid}

TaxaExpect receives habitat-annotated occurrence data from TaxaHabitat
and produces spatially-explicit priors for TaxaAssign (posterior
computation).

\textbf{Ecosystem:} TaxaFetch -\textgreater{} TaxaHabitat
-\textgreater{} \textbf{TaxaExpect} -\textgreater{} TaxaAssign

See the \href{https://github.com/DOI-USGS/TaxaID}{TaxaID README} for
ecosystem overview and installation instructions.

\subsection{Citation}\label{citation}

Lafferty, K.D., 2026, TaxaID -- A modular R ecosystem for Bayesian
taxonomic assignment: U.S. Geological Survey software release,
\url{https://doi.org/10.5066/xxxxxx}.

\subsection{Software Requirements}\label{software-requirements}

\begin{itemize}
\tightlist
\item
  R (\textgreater= 4.1.0; R Core Team 2025)
\item
  TaxaTools (foundation package)
\item
  TaxaFetch, TaxaHabitat (for fetching and habitat-labeling occurrence
  records upstream of the kernel pathway; in Suggests)
\end{itemize}

All dependencies are declared in the DESCRIPTION file and installed
automatically.

Developed with \href{https://claude.ai/code}{Claude Code} (Anthropic
PBC, San Francisco, California).

\subsection{References}\label{references}

Chao, A. (1984). Nonparametric estimation of the number of classes in a
population. \emph{Scandinavian Journal of Statistics}, 11(4), 265--270.

Good, I.J. (1953). The population frequencies of species and the
estimation of population parameters. \emph{Biometrika}, 40(3-4),
237--264.

Kish, L. (1965). \emph{Survey Sampling}. New York: Wiley.

R Core Team (2025). R: A Language and Environment for Statistical
Computing. V.4.5.2. R Foundation for Statistical Computing, Vienna,
Austria. \url{https://www.r-project.org}

\end{document}
